\chapter{Giới thiệu}

\section{Bối cảnh và Động lực Nghiên cứu}

Ngành bán lẻ tạp hóa và siêu thị hiện đại đang chuyển dịch nhanh chóng từ công cụ tìm kiếm thụ động và đặt hàng thủ công bằng bảng tính sang \textbf{AI tác nhân (\emph{Agentic AI})}---các phần mềm tác nhân có khả năng tiếp nhận đầu vào đa phương thức (giọng nói, văn bản, hình ảnh nhãn sản phẩm và chuỗi thời gian giao dịch), gọi công cụ có cấu trúc, suy luận dưới ràng buộc nghiệp vụ và phối hợp an toàn với con người. Khác với chatbot hỏi-đáp tĩnh hay các mô hình dự báo chạy tách rời, một \textbf{tác nhân AI (\emph{AI Agent})} vận hành trong vòng lặp \emph{quan sát -- suy luận -- hành động có kiểm soát}: phân giải ý định tự nhiên còn mơ hồ, truy vấn cơ sở dữ liệu sự thật hoặc động cơ toán học xác định, kiểm tra ranh giới an toàn/chính sách, rồi trả về khuyến nghị hoặc đề xuất hành động có thể kiểm chứng.

Trên thế giới, các tập đoàn bán lẻ hàng đầu giai đoạn 2024--2026 đã chứng minh rõ cả tiềm năng thương mại lẫn thách thức kỹ thuật khi đưa tác nhân AI vào vận hành chuỗi siêu thị:
\begin{itemize}
    \item \textbf{Walmart (\emph{Sparky}):} Triển khai tác nhân mua sắm tự chủ có khả năng biến mục tiêu tổng quát của khách hàng (lên thực đơn bữa ăn, chuẩn bị tiệc theo ngân sách) hoặc ảnh chụp danh sách viết tay thành giỏ hàng cụ thể, đề xuất thay thế sản phẩm và tự động hóa tác vụ thương mại, qua đó gia tăng rõ rệt giá trị giỏ hàng~\cite{walmart_sparky}.
    \item \textbf{Amazon (\emph{Rufus}):} Tích hợp tác nhân mua sắm tạo sinh được huấn luyện trên thuộc tính danh mục, đánh giá người dùng và dữ liệu hỏi-đáp để trả lời các câu hỏi so sánh phức tạp qua văn bản, giọng nói và hình ảnh, tiến tới hỗ trợ mua sắm bán tự động dựa trên quy mô hệ sinh thái khép kín~\cite{amazon_rufus}.
    \item \textbf{Carrefour (\emph{Hopla} / \emph{Hopla+}):} Tiên phong tại châu Âu trong tư vấn thực phẩm đàm thoại bằng cách gắn mục tiêu dinh dưỡng (giảm đường, tối ưu Nutri-Score, lên công thức món ăn theo ngân sách) trực tiếp vào danh mục siêu thị, đồng thời phát triển trợ lý AI nội bộ hỗ trợ vận hành cửa hàng~\cite{carrefour_hopla}.
\end{itemize}

Tại \textbf{Việt Nam}, thị trường bán lẻ hiện đại thuộc nhóm tăng trưởng và chuyển đổi số nhanh nhất Đông Nam Á, với các chuỗi lớn đã ứng dụng AI vào nhiều khâu cốt lõi:
\begin{itemize}
    \item \textbf{Thanh Lam Mart (Thanh Lam Mart / Thanh Lam Mart):} Với quy mô hơn 4.500--5.000 điểm bán trên toàn quốc, Thanh Lam Mart xem dữ liệu lớn và học máy là xương sống vận hành nhằm tự động hóa đặt hàng tập trung, giảm tải kiểm kê thủ công và giải phóng tới 30\% thời gian của nhân viên cửa hàng để tập trung chăm sóc khách hàng~\cite{vneconomy_retail_ai}.
    \item \textbf{Bách Hóa Xanh (MWG):} Cạnh tranh trực diện ở phân khúc thực phẩm tươi sống và tiêu dùng nhanh, Bách Hóa Xanh dựa vào phân phối định hướng dữ liệu và kiểm soát tồn kho ngắn hạn để cắt giảm hao hụt hàng dễ hỏng (\emph{shrinkage}) mà vẫn bảo đảm kệ hàng tươi mới~\cite{bhx_supplychain}.
    \item \textbf{FPT Retail (Long Châu, FPT Shop) \& FPT.AI:} Kết hợp tối ưu hóa chuỗi cung ứng bằng AI với tổng đài giọng nói tiếng Việt (\emph{Voicebot CSAT}) và chatbot kết nối ERP/CRM giúp tự động hóa tra cứu và chăm sóc khách hàng 24/7 ở quy mô lớn~\cite{fpt_ai_retail}.
\end{itemize}

Tuy nhiên, một khoảng trống lớn vẫn tồn tại giữa hệ thống phân tích ở cấp trụ sở với \textbf{thực thi ngay tại từng cửa hàng bán lẻ ở Việt Nam}. Tại \textbf{sàn bán hàng}, khách mua sắm trực tiếp hiếm khi được hỗ trợ bởi một trợ lý giọng nói tiếng Việt phản hồi tức thì, chịu nhiễu siêu thị tốt, so sánh chính xác đơn giá quy đổi, khuyến mãi, mức đường và lên giỏ món ăn theo ngân sách trên hàng nghìn thương hiệu Việt--quốc tế (\emph{Thanh Lam Home}, \emph{TH True Milk}, \emph{ACEFOODS}, \emph{OMO}) mà không bịa giá, bịa tồn kho hay đưa lời khuyên y tế rủi ro. Trong khi đó, tại \textbf{hậu cần cửa hàng}, cửa hàng trưởng tại các siêu thị mini và cửa hàng tiện lợi vẫn xoay xở thủ công với đơn đặt hàng mỗi ngày, lô thực phẩm tươi sống hạn dùng 1--3 ngày, biến động nhu cầu đột xuất và hàng trả nhà cung cấp---thiếu công cụ khôi phục nhu cầu bị che khuất khi đứt hàng và thiếu cơ chế phê duyệt Con người trong Vòng lặp (HITL) an toàn.

Đồ án tốt nghiệp này giải quyết đồng thời cả hai mặt trận trong một kiến trúc tác nhân kép thống nhất, triển khai qua hai kho mã nguồn cùng thuộc một dự án: \texttt{Retail-AI-} (tác nhân kiosk giọng nói tiếng Việt tên \emph{Mai} phục vụ khách hàng) và \texttt{ops\_chest} (\url{https://github.com/quanchu14104/ops_chest.git}, tác nhân vận hành hàng tươi sống phục vụ cửa hàng trưởng).

\section{Tổng quan Tài liệu và Hệ thống Liên quan}

Xây dựng tác nhân AI đáng tin cậy cho chuỗi bán lẻ đòi hỏi tổng hợp bốn nhánh nghiên cứu và kỹ nghệ then chốt: (i) tác nhân LLM gọi công cụ và ranh giới an toàn, (ii) khử nhiễu tiếng nói thời gian thực và ASR truyền luồng trong môi trường tạp âm, (iii) dự báo nhu cầu và tối ưu tồn kho hàng dễ hỏng theo lợi nhuận, và (iv) các hệ thống AI bán lẻ thực tế.

\textbf{Tác nhân AI dựa trên LLM, RAG có căn cứ và Cổng an toàn.}
Yao và cộng sự~\cite{yao2023react} đề xuất khung \emph{ReAct}, chứng minh việc đan xen suy luận với gọi công cụ bên ngoài giảm mạnh ảo giác so với sinh chuỗi suy nghĩ thuần túy. Schick và cộng sự~\cite{schick2023toolformer} cùng Lewis và cộng sự~\cite{lewis2020rag} khẳng định việc giao tác vụ tính toán, truy xuất và tra cứu sự thật cho công cụ xác định nâng cao độ chính xác thông tin. Dù vậy, vòng lặp tác nhân thiếu kiểm soát gây ra hai lỗi nghiêm trọng trong môi trường thực: \emph{khuếch đại độ trễ} (mỗi vòng LLM lập kế hoạch/duyệt câu trả lời cộng thêm 2--4\,giây chờ chết) và \emph{lệch chuẩn an toàn}. Röttger và cộng sự~\cite{rottger2024xstest} qua bộ kiểm thử \emph{XSTest} chỉ ra các bộ lọc an toàn thô sơ hay từ chối nhầm câu hỏi vô hại, còn OWASP Top 10 cho LLM~\cite{owasp_llm} cảnh báo tiêm lệnh (\emph{prompt injection}), rò rỉ dữ liệu và lạm quyền thực thi (\emph{excessive agency}). Vì vậy, tác nhân bán lẻ thực tế phải có cổng tiền-LLM xác định, kiểm chứng số liệu thời gian thực và tách biệt quyền tạo đề xuất khỏi quyền thực thi.

\textbf{Khử nhiễu Tiếng nói Thời gian thực và ASR Truyền luồng.}
Kiosk đặt tại siêu thị phải hoạt động giữa mức nhiễu 5--15\,dB gồm tiếng ù điều hòa/điện lưới 50/100\,Hz, tiếng xe đẩy và tiếng nói nền nhiều người. Các phương pháp lọc phổ cổ điển dễ làm méo thanh điệu tiếng Việt, trong khi các mạng khử nhiễu nơ-ron hiện đại như \emph{DeepFilterNet}~\cite{schroter2022deepfilternet} và hậu duệ gọn nhẹ \texttt{DPDFNet2}~\cite{sherpa_onnx} khử nhiễu không dừng trên CPU với méo tín hiệu rất thấp. Khi kết hợp với mô hình phát hiện tiếng nói \emph{Silero VAD}~\cite{silero_vad} và ASR truyền luồng đám mây hỗ trợ tăng trọng từ vựng miền (\emph{Azure AI Speech} \code{PhraseListGrammar}~\cite{azure_speech}), âm thanh được giải mã song song ngay khi khách đang nói thay vì đợi tải cả tệp WAV sau khi im lặng. Dù vậy, thực tế kiosk vẫn phải xử lý triệt để tiếng vọng loa khi khách ngắt lời và lỗi nhận diện tên thương hiệu tiếng Anh nói theo giọng Việt (như \emph{Thanh Lam Home} bị nghe thành \emph{Thanh Lam Mart hôn}).

\textbf{Dự báo Nhu cầu và Tối ưu Tồn kho Hàng Tươi sống.}
Nghiên cứu quản trị tồn kho hàng dễ hỏng chỉ ra hai bẫy kinh điển. Thứ nhất, dữ liệu bán hàng POS bị \emph{kiểm duyệt phải}: khi kệ đứt hàng từ giữa ngày, doanh số ghi nhận thấp hơn nhu cầu thực, đẩy mô hình dự báo vào vòng xoáy giảm đặt hàng liên tiếp~\cite{nahmias2011perishable}. Bộ dữ liệu mở \emph{FreshRetailNet-50K}~\cite{freshretailnet50k} cung cấp nhãn đứt hàng theo giờ tại các cửa hàng tiện lợi thực, cho phép khôi phục nhu cầu thực và huấn luyện mô hình dự báo phân vị bằng cây tăng cường độ dốc (\emph{LightGBM}~\cite{ke2017lightgbm}). Thứ hai, tối ưu sai số thống kê (như WAPE) không đồng nghĩa tối đa hóa lợi nhuận cửa hàng với hàng hạn dùng 1--3 ngày, vì chi phí dư hàng (hủy bỏ 100\% giá vốn) và thiếu hàng (mất lãi gộp và mất khách) bất đối xứng, chịu chi phối bởi phân vị tới hạn \emph{Newsvendor}~\cite{silver2017inventory} và cơ chế xuất bán theo lô cận hạn FEFO.

\textbf{Tổng hợp So sánh và Khoảng trống Nghiên cứu.}
Bảng~\ref{tab:lit_comparison} so sánh các phương pháp học thuật và hệ thống AI bán lẻ tiêu biểu với yêu cầu của chuỗi cửa hàng Việt Nam.

\begin{table}[H]
\centering
\caption{So sánh các phương pháp nghiên cứu và hệ thống AI bán lẻ tiêu biểu}
\label{tab:lit_comparison}
\footnotesize
\begin{tabularx}{\textwidth}{L{2.9cm} L{2.7cm} L{4.1cm} Y}
\toprule
\textbf{Công trình / Hệ thống} & \textbf{Cách tiếp cận} & \textbf{Ưu điểm chính} & \textbf{Hạn chế với bài toán tại Việt Nam} \\
\midrule
ReAct~\cite{yao2023react} \& RAG~\cite{lewis2020rag} & Suy luận LLM đa bước + gọi công cụ truy xuất & Linh hoạt với hội thoại mở nhiều lượt; giảm ảo giác tham số & Độ trễ cao (5--14\,s) do gọi LLM tuần tự; thiếu ranh giới đơn giá, an toàn và HITL bán lẻ \\
Walmart Sparky~\cite{walmart_sparky} \& Amazon Rufus~\cite{amazon_rufus} & Tác nhân mua sắm đa kênh độc quyền & Chuyển đổi giỏ hàng cao, hỗ trợ ảnh/danh sách, tích hợp đánh giá sâu & Hệ sinh thái đóng cho tiếng Anh/Mỹ; không giải bài toán giọng nói kiosk tiếng Việt nhiều nhiễu và vận hành cửa hàng \\
Carrefour Hopla~\cite{carrefour_hopla} & Chatbot thực phẩm dinh dưỡng + lên giỏ món ăn & Tư vấn giỏ hàng theo ngân sách và dinh dưỡng tốt & Thiên về web/văn bản; tách rời khỏi quản lý lô cận hạn FEFO và đặt hàng hậu cần \\
Thanh Lam Mart / BHX / FPT.AI~\cite{vneconomy_retail_ai,fpt_ai_retail} & Đặt hàng ML tập trung + voicebot tổng đài & Chứng minh hiệu quả quy mô chuỗi Việt Nam; giảm thời gian kiểm kê & Phân tích cấp trụ sở và tổng đài vẫn tách rời kiosk tư vấn giọng nói tại quầy và trợ lý HITL minh bạch cho cửa hàng trưởng \\
DeepFilterNet~\cite{schroter2022deepfilternet} \& Silero VAD~\cite{silero_vad} & Khử nhiễu nơ-ron cục bộ + phát hiện tiếng nói & Khử tạp âm trên CPU theo thời gian thực, bắt điểm dừng chính xác & Mô-đun âm thanh độc lập; cần tích hợp với STT truyền luồng, sửa lỗi thương hiệu và chống vọng ngắt lời \\
LightGBM~\cite{ke2017lightgbm} trên FreshRetailNet~\cite{freshretailnet50k} & Dự báo nhu cầu phân vị bằng cây tăng cường độ dốc & Học tốt độ trễ phi tuyến, chu kỳ thứ/lễ và lịch sử đứt hàng & Tối ưu WAPE đơn thuần có thể đặt dư và làm giảm lợi nhuận trên cửa hàng nhỏ nếu thiếu cổng định tuyến theo lợi nhuận \\
\bottomrule
\end{tabularx}
\end{table}

\textbf{Khoảng trống Nghiên cứu và Kỹ nghệ Mục tiêu.}
Như Bảng~\ref{tab:lit_comparison} chỉ ra, các hệ thống hiện hành mới giải quyết từng mảnh rời rạc. Khoảng trống cốt lõi mà đồ án hướng tới là thiết kế một \textbf{kiến trúc tác nhân AI kép độ trễ thấp, có kiểm chứng, may đo cho chuỗi bán lẻ Việt Nam}: (1) cắt bỏ hoàn toàn các vòng planner LLM thừa tại quầy bán hàng mà vẫn giữ 100\% căn cứ danh mục và an toàn y tế/OWASP; (2) duy trì nhận dạng giọng nói tiếng Việt kèm thương hiệu tiếng Anh chính xác dưới nhiễu siêu thị 12\,dB--5\,dB với độ trễ sau nói dưới 1 giây; và (3) kết hợp dự báo hàng tươi sống định tuyến theo cổng lợi nhuận (LightGBM/EWMA) và mô phỏng lô FEFO với kiểm soát số liệu thời gian thực và phê duyệt HITL bắt buộc.

\section{Phát biểu Bài toán}

Tại cấp độ cửa hàng trong một chuỗi siêu thị hoặc cửa hàng tiện lợi tại Việt Nam, hai nhóm vấn đề tương hỗ đồng thời làm suy giảm trải nghiệm khách hàng và biên lợi nhuận vận hành:

\begin{enumerate}
    \item \textbf{Bài toán Tư vấn Khách hàng tại Sàn bán hàng (\texttt{Retail-AI-}):} Khách hàng đặt câu hỏi tự nhiên bằng tiếng Việt trong không gian ồn ào: lọc thương hiệu (\emph{``khăn mặt cotton Thanh Lam Home''}), nhu cầu dinh dưỡng (\emph{``sữa tươi không đường''}), so sánh tiết kiệm theo đơn giá (\emph{``nước giặt OMO loại nào tiết kiệm nhất?''}), so sánh hai hãng hoặc lên giỏ nấu ăn theo ngân sách (\emph{``mua đồ nấu lẩu thái dưới 300k''}). Kiến trúc LLM thông thường vấp phải 4 rào cản: (a) \emph{độ trễ đa lượt planner} (2--4 lượt gọi LLM tuần tự không streaming gây chờ 5,4--14,2\,s); (b) \emph{suy giảm nhận dạng giọng nói trong tạp âm và trễ tải tệp} (tải WAV sau khi nói lên Whisper đám mây tốn 2,0--3,0\,s và WER 14,2--22,2\% trong nhiễu, trong khi sửa lỗi từ vựng thô sơ làm hỏng từ tiếng Việt thông thường); (c) \emph{tiếng vọng loa bán song công} cản trở ngắt lời tự nhiên; và (d) \emph{rủi ro ảo giác và an toàn} về giá bán, tồn kho, dị ứng hay tư vấn thuốc.
    \item \textbf{Bài toán Vận hành Hậu cần Hàng Tươi sống (\texttt{ops\_chest}):} Cửa hàng trưởng quản lý hàng trăm SKU tươi sống hạn dùng 1--3 ngày phải quyết định mỗi ngày số lượng nhập, mức giảm giá cận hạn, lô trả nhà cung cấp và lô hủy. Cách đặt hàng theo kinh nghiệm chịu 14,5\% đứt hàng chung, 8,7\% đứt hàng nhóm hút khách và 5,3\% hủy hàng, trong khi doanh số thô bị thiếu hụt $-32,3\%$ ở các ngày đứt hàng. Ngược lại, áp dụng ML dự báo mù quáng cho mọi phân khúc lại đặt dư trên cửa hàng nhỏ (giảm $-5,6\%$ lợi nhuận trung vị), còn cấp quyền tự thực thi cho chatbot LLM tạo ra rủi ro vận hành nghiêm trọng nếu mô hình bịa số liệu hoặc tự nhận đã gửi đơn đặt hàng khi chưa được duyệt.
\end{enumerate}

Câu hỏi nghiên cứu trung tâm là: \emph{Làm thế nào để xây dựng một kiến trúc tác nhân AI kép thống nhất vừa giải quyết bài toán tư vấn giọng nói tiếng Việt tại quầy theo thời gian thực không ảo giác, vừa trang bị cho cửa hàng trưởng trí tuệ tồn kho hàng tươi sống được kiểm chứng theo lợi nhuận dưới cơ chế quản trị Con người trong Vòng lặp (HITL) nghiêm ngặt?}

\section{Mục tiêu và Đóng góp của Đồ án}

\subsection*{Mục tiêu Tổng quát}
Thiết kế, hiện thực hóa và đánh giá thực nghiệm một nền tảng tác nhân AI kép cho chuỗi cửa hàng bán lẻ Việt Nam, kết hợp trợ lý kiosk giọng nói chịu nhiễu thời gian thực (\texttt{Retail-AI-}) và trợ lý vận hành hàng tươi sống định tuyến theo lợi nhuận có phê duyệt HITL (\texttt{ops\_chest}).

\subsection*{Các Mục tiêu Cụ thể}
\begin{itemize}
    \item \textbf{Xây dựng chuỗi xử lý giọng nói tiếng Việt độ trễ thấp, chịu nhiễu siêu thị (\texttt{Retail-AI-}):} Tích hợp Silero VAD~v6 cục bộ, STT truyền luồng WebSocket thời gian thực (\code{/audio/transcribe/stream}) chạy Azure AI Speech SDK kèm \code{PhraseListGrammar} Thanh Lam Mart, mô hình khử nhiễu \texttt{DPDFNet2} ONNX thích nghi SNR (bỏ qua mic sạch và cứu hộ clip nhiễu nặng), sửa lỗi thương hiệu neo theo danh mục, ngắt lời song công chống vọng 3 lớp và Azure Neural TTS theo mệnh đề.
    \item \textbf{Loại bỏ độ trễ planner thừa bằng phát luồng SSE có căn cứ (\texttt{Retail-AI-}):} Thiết kế động cơ streaming không planner (\code{streaming_consultation.py}) kết hợp bộ lắp giỏ hàng 1 lượt gọi LLM giúp trả thẻ sản phẩm xác thực trong vài chục mili-giây và phát luồng token DeepSeek trực tiếp, đồng thời định tuyến các câu hỏi y tế, dị ứng, trẻ nhỏ và tấn công OWASP qua cổng an toàn xác định dưới 3\,ms.
    \item \textbf{Triển khai dự báo định tuyến lợi nhuận và vận hành HITL (\texttt{ops\_chest}):} Hiện thực hóa khôi phục nhu cầu đứt hàng theo giờ, dự báo phân vị LightGBM~v3 (29 đặc trưng, P10/P50/P90), bộ định tuyến cổng lợi nhuận ML/EWMA kiểm định chéo, mô phỏng nhập hàng Newsvendor có trần hạn dùng FEFO, quy trình xử lý lô cận hạn và tác nhân 31 công cụ kèm kiểm soát số liệu thời gian thực và phê duyệt quản lý bắt buộc.
    \item \textbf{Chuẩn hóa giao ước tích hợp liên phân hệ và kiểm chứng định lượng toàn diện:} Xác định giao ước API chỉ đọc kết nối \texttt{ops\_chest} với \texttt{Retail-AI-}, đồng thời kiểm chứng hai kho mã nguồn qua bộ test tự động, 100 câu hỏi SOTA~V2, thực nghiệm 5 cấu hình STT trong nhiễu, 10 truy vấn streaming thực tế và backtest lợi nhuận trên 50 cửa hàng độc lập.
\end{itemize}

\subsection*{Đóng góp Chính}
\begin{itemize}
    \item \textbf{Đóng góp về phương pháp luận và kiến trúc:} Chứng minh ba nguyên lý thiết kế tác nhân AI có tính tổng quát cao: (i) \emph{phát luồng SSE không planner kết hợp định tuyến lai Groq LPU + bộ đệm ngữ nghĩa Turn-1 và biên khu vực Singapore (\texttt{sin1})}---tách biệt văn xuôi phát luồng tốc độ cao khỏi lập kế hoạch JSON an toàn đa bước giúp hạ độ trễ HTTPS thực tế từ 2,4--5,5\,s xuống \textbf{508--870\,ms} lạnh và \textbf{125--140\,ms} qua bộ đệm (0 token) mà không đánh đổi độ chính xác dữ liệu; (ii) \emph{làm giàu tri thức sản phẩm đa nguồn 3 tầng bằng chứng kết hợp pipeline làm sạch 8 tầng đạt 0 lỗi} (theo chuẩn Hugging Face \code{datatrove} và OpenFoodFacts) giúp mở rộng độ phủ bảng thành phần chi tiết lên 1.214+18 SKU mà vẫn fail-closed trung thực với hàng chưa có nhãn; và (iii) \emph{định tuyến dự báo theo cổng lợi nhuận}---chọn ML hay thống kê dựa trên lợi nhuận mô phỏng ngoài mẫu và ghép trung vị LightGBM với độ lệch chuẩn EWMA vượt trội hoàn toàn cách chọn theo WAPE thuần túy.
    \item \textbf{Đóng góp thực tiễn:} Cung cấp hai nguyên mẫu phần mềm hoàn chỉnh, được kiểm thử sâu cho bán lẻ Việt Nam: kiosk giọng nói và nhận diện ảnh di động tiếng Việt đã triển khai trực tuyến trên Vercel Singapore (\url{https://mai-retail-kiosk.vercel.app}) đạt 0,0\% WER thương hiệu ở nhiễu 12\,dB, popup xem ảnh sản phẩm phủ 99,9\% danh mục, báo quản lý thời gian thực qua Zalo Bot; và trợ lý vận hành đóng gói Docker hiệu chuẩn trên dữ liệu tươi sống thực giúp tăng $+4,9\%$ lợi nhuận gộp cửa hàng trong khi bảo toàn 100\% trách nhiệm phê duyệt của con người.
\end{itemize}

\section{Phạm vi và Giới hạn}

\subsection*{Phạm vi Triển khai}
Dự án hoàn thiện và kiểm chứng cả hai phân hệ ở mức nguyên mẫu kỹ thuật tiền thương mại: \texttt{Retail-AI-} vận hành song song cục bộ (\code{127.0.0.1:8000} / \code{5173}) và triển khai trực tuyến trên Vercel Production khu vực Singapore (\texttt{sin1}, \url{https://mai-retail-kiosk.vercel.app}) trên danh mục 3.099 SKU Thanh Lam Mart (kèm 2.137 SKU đã làm giàu tri thức thành phần/cách dùng/bảo quản), giao diện di động QR có popup ảnh sản phẩm, cảnh báo quản lý qua Zalo Bot và giỏ hàng/thanh toán mô phỏng; \texttt{ops\_chest} chạy trên ngăn xếp Docker FastAPI + React/Vite + PostgreSQL + Mailpit được hiệu chuẩn trên nhu cầu thực FreshRetailNet-50K (Cửa hàng mẫu 18 và 50 cửa hàng kiểm định) cùng môi trường dữ liệu tổng hợp 263 SKU cho kiểm thử quy tắc nghiệp vụ.

\subsection*{Giới hạn Hiện tại}
\begin{itemize}
    \item \textbf{Dữ liệu danh mục chụp tại một thời điểm và giao dịch mô phỏng:} \texttt{Retail-AI-} sử dụng bản chụp danh mục Thanh Lam Mart ngày 15/09/2026 kèm dữ liệu lượt bán 30 ngày gắn nhãn \texttt{SIMULATED} tách biệt và thanh toán sandbox, chưa kết nối trực tiếp POS/ERP thời gian thực của chuỗi siêu thị.
    \item \textbf{Chưa đồng bộ trực tuyến thời gian thực giữa hai kho mã nguồn:} Dù thuộc chung một kiến trúc thống nhất với giao ước chỉ đọc rõ ràng (Phụ lục~C), việc đồng bộ trực tiếp sổ lô tồn kho từ \texttt{ops\_chest} sang danh mục kiosk của \texttt{Retail-AI-} mới ở mức đặc tả giao diện, chưa triển khai trên hạ tầng chuỗi cửa hàng thật.
    \item \textbf{Khoảng cách phòng thí nghiệm--siêu thị thật và độ phủ nhãn dinh dưỡng/calo:} Dù kiến trúc lai Groq LPU + Semantic Cache + vùng \texttt{sin1} đã khắc phục ngưỡng trễ 1,2--1,5\,s của nhánh tư vấn phát luồng, các ca JSON có cấu trúc phức tạp (dị ứng/y tế sâu) vẫn phụ thuộc API đám mây và khoảng 59\% SKU chưa có bảng thành phần định lượng hoặc calo công khai (đều được hệ thống báo trung thực \texttt{UNKNOWN}); thực nghiệm chịu nhiễu sử dụng hỗn hợp tạp âm 12\,dB và 5\,dB có kiểm soát thay vì thử nghiệm thực địa nhiều tuần tại siêu thị vật lý.
    \item \textbf{Tham số thương mại giả định trong \texttt{ops\_chest}:} Trong khi chuỗi nhu cầu hằng ngày của FreshRetailNet-50K là dữ liệu thực, tên hiển thị SKU, tỷ lệ giá vốn (62--75\%), hạn dùng 1--3 ngày, lịch giao hàng và hộp thư Mailpit là các giả định vận hành có tài liệu hóa minh bạch ở Chương~3.
\end{itemize}
